%% Théorème d'Al-Kashi dans un triangle quelconque ABC.
%% Convention du manuel : a = BC (côté opposé à A), b = AC, c = AB.
%% Le triangle n'est pas rectangle : Pythagore ne s'applique pas, Al-Kashi le généralise.
\documentclass[tikz,border=4pt]{standalone}
\usepackage{liaisons}
\begin{document}
\begin{tikzpicture}[x=1.4cm,y=1.4cm,
  cote/.style={line width=1.1pt, line join=round},
  etiq/.style={font=\small, inner sep=3pt}]

\coordinate (A) at (0,0);
\coordinate (B) at (4,0);
\coordinate (C) at (1.2,2.4);

%% ---------------- les trois côtés, étiquetés par leur longueur -----------
\draw[cote] (A) -- (B) node[midway, below, etiq] {$c=AB$};
\draw[cote] (A) -- (C) node[midway, sloped, above, etiq] {$b=AC$};
\draw[cote] (C) -- (B) node[midway, sloped, above, etiq] {$a=BC$};

%% ---------------- angle  au sommet A ------------------------------------
\draw[solideA, line width=1pt] (0.6,0) arc (0:63.435:0.6);
\node[font=\small, text=solideA, anchor=west, inner sep=2pt] at (31:0.72) {$\hat{A}$};

%% ---------------- sommets --------------------------------------------------
\fill (A) circle (1.8pt) node[etiq, anchor=north east] {$A$};
\fill (B) circle (1.8pt) node[etiq, anchor=north west] {$B$};
\fill (C) circle (1.8pt) node[etiq, anchor=south] {$C$};

%% ---------------- le théorème ----------------------------------------------
\rappel{\node[font=\small, anchor=north, inner sep=2pt] at (2,-0.62)
     {$a^{2}=b^{2}+c^{2}-2bc\cos(\hat{A})$};}
\end{tikzpicture}
\end{document}
